\abstract{
%---------- Motivation & Gaps
Ensemble datasets are ever more prevalent in various scientific domains.
In climate science, ensemble datasets are used to capture variability in projections under plausible future conditions including greenhouse and aerosol emissions.
Each ensemble model run produces projections that are fundamentally similar yet meaningfully distinct.
Understanding this variability among ensemble model runs and analyzing its magnitude and patterns is a vital task for climate scientists.
%---------- Our Approach
In this paper, we present \ClimateSOM, a visual analysis workflow that leverages a self-organizing map (SOM) and Large Language Models (LLMs) to support interactive exploration and interpretation of climate ensemble datasets.
The workflow abstracts climate ensemble model runs---spatiotemporal time series---into a distribution over a 2D space that captures the variability among the ensemble model runs using a SOM.
LLMs are integrated to assist in sensemaking of this SOM-defined 2D space, the basis for the visual analysis tasks.
In all, \ClimateSOM enables users to explore the variability among ensemble model runs, identify patterns, compare and cluster the ensemble model runs.
%---------- Results
To demonstrate the utility of \ClimateSOM, we apply the workflow to an ensemble dataset of precipitation projections over California and the Northwestern United States.
Furthermore, we conduct an expert review of the visual workflow and the insights from the case studies with six domain experts to evaluate our approach and its utility.
}


